% Source: https://artofproblemsolving.com/community/c3h2604632_wlog_question
% Stable identity: aops-topic-2604632-q01
Factor the right-hand side: $ab+bc-ca-b^2=b(a-b)-c(a-b)=(a-b)(b-c)$. So the conditions are
\[a+b+c=20,\qquad(a-b)(b-c)=20.\]

\textbf{(b)} The conditions are unchanged if $a$ and $c$ are swapped, since $(c-b)(b-a)=(a-b)(b-c)$. Also $(a-b)(b-c)=20>0$, so $a-b$ and $b-c$ have the same sign: either $a>b>c$ or $a<b<c$. Swapping $a$ and $c$ turns the second case into the first. So assuming $a>b>c$ is justified, but only because of this argument. The general assumption ``$a\ge b\ge c$'' would not be allowed for an arbitrary non-symmetric equation; for example, swapping $a$ and $b$ changes this one.

\textbf{(a)} Take $a>b>c$ and put $p=a-b$, $q=b-c$, positive integers with $pq=20$. Then $a+b+c=3b+p-q=20$, so $b=\dfrac{20-p+q}{3}$ and $c=b-q$.
\begin{itemize}
\item $(p,q)=(1,20)$: $b=13$, $c=-7$, impossible.
\item $(p,q)=(2,10)$: $b=\dfrac{28}{3}$, not an integer.
\item $(p,q)=(4,5)$: $b=7$, $a=11$, $c=2$.
\item $(p,q)=(5,4)$: $b=\dfrac{19}{3}$, not an integer.
\item $(p,q)=(10,2)$: $b=4$, $a=14$, $c=2$.
\item $(p,q)=(20,1)$: $b=\dfrac13$, not an integer.
\end{itemize}
So $(a,b,c)$ is $(11,7,2)$ or $(14,4,2)$, or one of these with $a$ and $c$ swapped. Check: $(11-7)(7-2)=20$ and $(14-4)(4-2)=20$. The possible products are
\[abc=\boxed{154\text{ or }112}.\]
